% Part: methods
% Chapter: proofs
% Section: using-definitions

\documentclass[../../../include/open-logic-section]{subfiles}

\begin{document}

\olfileid{mod}{prf}{def}

\olsection{વ્યાખ્યાઓનો ઉપયોગ}

અગાઉ આપણે કહ્યું હતું કે સાબિતીમાં વપરાઈ શકે એવી બધી વ્યાખ્યાઓથી
તમારે પરિચિત હોવું જોઈએ અને તેમને યોગ્ય રીતે લાગુ પાડતા આવડવું
જોઈએ. આ ખૂબ અગત્યની વાત છે; તેથી તેને થોડું વધુ વિગતે જોઈએ.
વ્યાખ્યાઓ ગુણધર્મો અને સંબંધોનાં ટૂંકાં નામ આપે છે, જેથી આપણે
તેમના વિશે સંક્ષેપમાં વાત કરી શકીએ. રજૂ થયેલા ટૂંકા નામને
\emph{વ્યાખ્યેય પદ} (definiendum) અને તે જેના માટે વપરાય છે તેને
\emph{વ્યાખ્યાકારક વિધાન} (definiens) કહીએ. સાબિતીમાં આપણે ઘણી વાર
વ્યાખ્યેય પદ કેવી રીતે રજૂ થયું હતું તે તરફ પાછા જવું પડે છે:
સાબિતીને આગળ વધારવા વ્યાખ્યાકારક વિધાનની તાર્કિક રચનાનો ઉપયોગ
કરવો પડે છે. વ્યાખ્યાઓને ઉકેલતાં, તમે તર્કની મુખ્ય કામગીરી ક્યાં
થાય છે ત્યાં પહોંચો છો.

એક ઉદાહરણથી શરૂ કરીએ. ધારો કે તમારે નીચેનું સાબિત કરવું છે:

\begin{prop}
કોઈ પણ ગણો $A$ અને $B$ માટે $A \cup B = B \cup A$.
\end{prop}

સાબિતી શરૂ કરવા માટે પણ બે ગણો એકસમાન હોવાનો અર્થ જાણવો પડે;
એટલે, ઉપરના સમીકરણમાં ગણો માટે ``$=$''નો અર્થ શું છે તે જાણવું પડે.
ગણો ત્યારે જ એકસમાન છે જ્યારે તેમના !!{element}s સરખા હોય.
તેથી આપણે નીચેની વ્યાખ્યા ઉકેલવી પડે છે:

\begin{defn}
ગણો $A$ અને $B$ \emph{એકસમાન} છે, $A = B$, ત્યારે અને ત્યારે જ
જ્યારે~$A$નું દરેક !!{element}~$B$નું !!a{element} હોય અને ઊલટું પણ હોય.
\end{defn}

આ વ્યાખ્યામાં $A$ અને~$B$ મનસ્વી ગણો માટેનાં સ્થાનાપન્ન ચિહ્નો છે.
તે જેની વ્યાખ્યા આપે છે---\emph{વ્યાખ્યેય પદ}---તે ``$A = B$'' છે;
અને $A = B$ ક્યારે સાચું છે તેની શરત આપે છે. આ શરત---``~$A$નું દરેક
!!{element}~$B$નું !!a{element} છે અને ઊલટું પણ છે''---એ
\emph{વ્યાખ્યાકારક વિધાન} છે.\footnote{આ ખાસ કિસ્સામાં---અને આ
  ગૂંચવણ ઊભી કરે છે!---જ્યારે $A = B$ હોય, ત્યારે $A$ અને $B$
  ખરેખર એક જ ગણ છે, ભલે તેની ડાબી અને જમણી બાજુ માટે જુદા
  અક્ષરો વાપરીએ. પરંતુ તે ગણને દર્શાવવાની રીતો જુદી હોઈ શકે,
  અને તેથી આ વ્યાખ્યા તુચ્છ નથી.} વ્યાખ્યા કહે છે કે શરત પૂરી
થાય ત્યારે અને ત્યારે જ $A = B$ સાચું છે (અંગ્રેજીમાં આને
ટૂંકમાં ``iff'' લખાય છે).

વ્યાખ્યા લાગુ પાડતી વખતે તેમાંના $A$ અને $B$ને તમે જે ખાસ
કિસ્સો તપાસો છો તેની સાથે મેળવો. આપણા કિસ્સામાં,
$A \cup B = B \cup A$ સાચું હોવા માટે દરેક $z \in A \cup B$,
$B \cup A$માં પણ હોવું જોઈએ અને ઊલટું પણ. પ્રતિજ્ઞાપનમાંનું
$A \cup B$ વ્યાખ્યાના~$A$નું સ્થાન લે છે અને $B \cup A$ તેના~$B$નું.
કારણ કે $A$ અને $B$ બંને વ્યાખ્યા અને સાબિત કરવાના પ્રતિજ્ઞાપનમાં
વપરાય છે, પણ તેમની ભૂમિકાઓ જુદી છે, તેથી તેમને ગૂંચવી ન દેવા
ધ્યાન રાખો. દાખલા તરીકે,~$A$નું દરેક !!{element}~$B$નું
!!a{element} છે અને ઊલટું પણ છે તે બતાવીને પ્રતિજ્ઞાપન સાબિત
થઈ જશે એમ માનવું ભૂલ હશે: તે $A = B$ બતાવશે, $A \cup B = B \cup A$
નહીં. (વળી, $A$ અને $B$ કોઈ પણ બે ગણો હોઈ શકે; $A$ અને~$B$
વિશે કંઈ ધાર્યું ન હોય તો તેઓ જુદા ગણો પણ હોઈ શકે છે.)

સાબિતીમાં આપણે ગણસિદ્ધાંતના યોગગણ જેવા ખ્યાલો સાથે કામ કરીએ
છીએ; એટલે સાબિતી કેવી રીતે આગળ વધારવી તે સમજવા $\cup$ સંકેતનો
અર્થ પણ જાણવો પડે. ક્યારેક એક વ્યાખ્યા ઉકેલવાથી બીજી વ્યાખ્યાઓ
પણ ઉકેલવી પડે છે. દાખલા તરીકે, $A \cup B$ની વ્યાખ્યા
$\Setabs{z}{z \in A \text{ અથવા } z \in B}$ છે. તેથી $x \in A \cup
B$ સાબિત કરવું હોય, તો $\cup$ની વ્યાખ્યા ઉકેલતાં
$x \in \Setabs{z}{z \in A \text{ અથવા } z \in B}$ સાબિત કરવું પડે.
હવે એ પણ યાદ રાખો કે $x \in \Setabs{z}{\dots z\dots}$ ત્યારે અને
ત્યારે જ જ્યારે $\dots x\dots$. તેથી $\Setabs{z}{\dots z \dots}$
સંકેતની વ્યાખ્યા વધુ ઉકેલતાં તમારે $x \in A$ અથવા $x \in B$
બતાવવું પડે. એટલે ``$A \cup B$નું દરેક !!{element}, $B \cup A$નું
પણ !!a{element} છે'' તેનો ખરેખર અર્થ છે: ``દરેક $x$ માટે,
જો $x \in A$ અથવા $x \in B$, તો $x \in B$ અથવા $x \in A$.''
પ્રતિજ્ઞાપનમાંની વ્યાખ્યાઓ પૂરેપૂરી ઉકેલીએ તો આપણે નીચેનું
બતાવવાનું રહે છે:

\begin{prop}
કોઈ પણ ગણો $A$ અને $B$ માટે: (a) દરેક $x$ માટે, જો $x \in A$
અથવા $x \in B$, તો $x \in B$ અથવા $x \in A$; અને (b) દરેક
$x$ માટે, જો $x \in B$ અથવા $x \in A$, તો $x \in A$ અથવા
$x \in B$.
\end{prop}

અગત્યની વાત એ છે કે વ્યાખ્યાઓ ઉકેલવી એ સાબિતી રચવાનો આવશ્યક
ભાગ છે. યોગ્ય રીતે આમ કરવું ક્યારેક મુશ્કેલ હોય છે: વ્યાખ્યામાંના
ચલો અને તમે સાબિત કરતા દાવામાંના પદોને જુદાં ઓળખીને તેમની વચ્ચે
યોગ્ય મેળ બેસાડવો પડે. સફળ થવા માટે પ્રશ્ન શું પૂછે છે અને તેમાં
વપરાયેલા બધા પદોનો અર્થ શું છે તે જાણવું જરૂરી છે---ઘણી વાર
એકથી વધુ વ્યાખ્યાઓ ઉકેલવી પડે છે. નીચેનાં સરળ ઉદાહરણોમાં
ઉકેલ લગભગ તરત જ વ્યાખ્યાઓમાંથી નીકળે છે. અલબત્ત, હંમેશાં
આટલું સહેલું નહીં હોય.

\begin{prob}
ધારો કે તમારે $A \cap B \neq \emptyset$ સાબિત કરવું છે. અહીં
આવતી બધી વ્યાખ્યાઓ ઉકેલો, એટલે કે ``$\cap$'', ``='', કે
``$\emptyset$''નો ઉલ્લેખ કર્યા વિના આ વાત ફરીથી કહો.
\end{prob}

\end{document}
